\chapter{Contest: Point-Level Instance Arbitration}
\label{ch:contest}

% PRESUPUESTO: ~8 pp. El capítulo más largo: es la contribución estrella.
% CONTRIBUCIÓN 4.
%
% ORDEN: este capítulo y el \ref{ch:geometric_consistency} son intercambiables.
% NO escribir referencias de orden ("como vimos en el capítulo anterior").

\section{From instance to point}
\label{sec:contest_granularity}
% Por qué un fallo PARCIAL exige una decisión PARCIAL. La decisión a nivel de instancia
% es todo-o-nada, y por eso no alcanza los fallos de \ref{sec:fm_scope_limit}.
% El punto es la unidad de verdad.

Como ya se ha visto, OVO separa dos tareas: la reconstrucción \emph{geométrica} del
mapa, que corre el SLAM, y una capa \emph{semántica} de asignación y clasificación que se
monta encima de esa reconstrucción. Cada \foreignlanguage{english}{keyframe} nuevo aporta
geometría, pero la reconstrucción no acumula todo lo que ve sin criterio: para mantener la
coherencia del mapa, antes de añadir puntos comprueba qué parte de lo observado
corresponde a superficie ya reconstruida y la descarta. Solo las regiones vistas por
primera vez generan puntos nuevos; lo ya reconstruido no se duplica, únicamente anota que
se ha vuelto a observar. El resultado es una nube \emph{deduplicada}: un punto por
superficie física, no uno por observación.

Es una decisión de diseño sana, y por dos motivos. Acota la escala del mapa al área de la
escena en vez de a la longitud de la trayectoria: una pared vista cien veces sigue siendo
una pared, no cien capas de puntos. Y evita apilar reconstrucciones ligeramente
desalineadas de la misma superficie, que la emborronarían. Esa permanencia geométrica está
bien ganada: la posición de una superficie es un \emph{hecho} que no cambia, y tratarla
como fija es lo correcto.

Sobre esa reconstrucción se monta la capa semántica. Cada punto lleva un dueño, la
instancia a la que pertenece, y ese dueño lo fija la primera máscara de SAM que reunió
votos suficientes para reclamar el punto, ya sea creando una instancia nueva o adjudicándolo
a una que ya existía y bajo la que cayó. La capa semántica hereda así la permanencia de la
geometría: una vez asignado, el punto conserva su dueño durante toda la vida del mapa. Pero las dos permanencias no valen lo mismo. La
posición de un punto es un hecho medido; su dueño es una hipótesis, la de una sola máscara
en un solo instante, que nunca se revisa.

De ahí la rigidez. Por muchos \foreignlanguage{english}{keyframes} que después muestren que
el dueño correcto es otro, no hay palanca para cambiarlo. El dedup lo agrava: la superficie
ya existe, así que la máscara que la reclamaría más tarde reproyecta sobre puntos
\emph{ajenos} y no crea puntos propios con los que competir. La disputa nunca aflora como
geometría; queda sepultada en el instante de la primera asignación.

Por eso el mecanismo de fusión de \foreignlanguage{english}{loop closure}
(\ref{sec:fm_scope_limit}) no puede alcanzar estos fallos. Ese \emph{merge} une dos
instancias cuando sus nubes se solapan, y el dedup impide que ese solape exista: los puntos
mal adjudicados no tienen copia correcta encima, son los mismos puntos con el dueño
equivocado. Sin dos nubes que comparar, el merge no tiene de qué dispararse. Tres
situaciones, recurrentes en la práctica, lo ejemplifican (Fig.~\ref{fig:contest_cases}).

\begin{figure}[htbp]
  \centering
  % PLACEHOLDER: sustituir cada \fbox por \includegraphics de un screenshot de Rerun (PNG),
  % cada uno un montaje Antes | Después sobre nube coloreada por instancia.
  % Escena + IDs de instancia por fijar (candidato para reusar como figura de fallo del Ch1).
  \newcommand{\caseph}[1]{\fbox{\parbox[c][2.2cm][c]{0.96\linewidth}{\centering\footnotesize\itshape #1}}}
  \begin{subfigure}{\linewidth}
    \centering\caseph{Antes: trozo contaminado \quad$\vert$\quad Después: trozo devuelto}
    \caption{Nacimiento bajo una máscara ajena.}
    \label{fig:contest_case_a}
  \end{subfigure}\\[4pt]
  \begin{subfigure}{\linewidth}
    \centering\caseph{Antes: revisita partida en dos instancias \quad$\vert$\quad Después: reunidas}
    \caption{Revisita parcial.}
    \label{fig:contest_case_b}
  \end{subfigure}\\[4pt]
  \begin{subfigure}{\linewidth}
    \centering\caseph{Antes: fragmento huérfano en el suelo \quad$\vert$\quad Después: reabsorbido}
    \caption{Granularidad inconsistente del segmentador.}
    \label{fig:contest_case_c}
  \end{subfigure}
  \caption{Los tres modos de fallo fuera del alcance del \emph{merge} (izquierda), resueltos
  por el contest (derecha).}
  \label{fig:contest_cases}
\end{figure}

\paragraph{Caso A: nacimiento bajo una máscara ajena.}
Un objeto entra en el campo de visión y SAM lo segmenta junto a otro que no es el suyo:
una sola máscara cubre ambos. Los puntos de los dos objetos nacen bajo la misma
instancia. Más tarde, cuando el objeto se ve despejado y SAM lo segmenta bien, la nueva
instancia se crea correctamente, pero el trozo que quedó atrapado en la primera ya tiene
dueño y no se lo devuelve nadie. La contaminación es permanente (Fig.~\ref{fig:contest_case_a}).

\paragraph{Caso B: la revisita parcial.}
Un objeto se ve parcialmente y desaparece de la escena antes de consolidarse. Al reentrar
por otro lado, la vista es mejor y se le asigna una instancia nueva que crece sana. Aunque
en algún \foreignlanguage{english}{keyframe} posterior las máscaras de SAM segmenten bien
el trozo que ya se había visto la primera vez, esos puntos siguen perteneciendo a la
instancia antigua. Las dos deberían ser una, pero no se solapan (son la misma superficie
vista en dos momentos, no dos copias en el mismo sitio), así que el merge no las
considera candidatas y jamás las une (Fig.~\ref{fig:contest_case_b}).

\paragraph{Caso C: granularidad inconsistente del segmentador.}
SAM no segmenta con una granularidad estable: a veces hila muy fino y parte una superficie
grande (el suelo es el caso típico), otras la da entera. Si el primer avistamiento de una
zona la parte en exceso y el resto de la trayectoria la ve completa, el fragmento inicial
queda como una instancia diminuta \emph{dentro} de la grande. Cae dentro de ella en todos
los frames siguientes, pero al tener ya dueño no se reabsorbe: el merge, de nuevo, no ve
duplicados que unir (Fig.~\ref{fig:contest_case_c}).

Los tres comparten raíz, y no es una heurística mal calibrada: son fallos \emph{fuera del
dominio} de los mecanismos existentes. Cada uno cumple su responsabilidad, el dedup la
coherencia geométrica y la fusión de \foreignlanguage{english}{loop closure} las
duplicidades por deriva o revisita, pero ninguno \emph{refina} las asignaciones a la luz de
lo que se observa después. Ese es el hueco que llena el contest, y no pasa por deshacer el
dedup sino por hacer revisable el dueño de un punto por evidencia en vez de fijo por la
primera máscara: la geometría sigue congelada, la semántica deja de estarlo. El
refinamiento descansa sobre dos decisiones.

\begin{enumerate}
  \item \textbf{Granularidad de la decisión.}
  Un fallo parcial solo se arregla con una decisión parcial. Si la unidad de decisión es la
  instancia entera, «este trozo es de B» no se puede ni plantear: o se mueve toda la instancia
  o ninguna. Bajando la decisión al punto, cada punto puede cambiar de dueño por su cuenta, y
  el \emph{merge} clásico pasa a ser el caso extremo en que resulta que todos los puntos de
  una instancia migran a otra.

  \item \textbf{Criterio para revocar.}
  Precisamente porque la primera decisión puede estar mal, revocarla a la ligera, con la
  evidencia de un solo frame, reintroduce el mismo problema en sentido contrario. La única
  base fiable para mover un punto de dueño es la \emph{evidencia acumulada}: observar a lo
  largo de muchos \foreignlanguage{english}{keyframes} qué instancia reclama ese punto de
  forma consistente, y solo entonces decidir. De confiar ciegamente en cada máscara se pasa a
  confiar en el historial de máscaras.
\end{enumerate}

\medskip
\noindent\emph{Precondición.} Todo esto presupone geometría estable y
\foreignlanguage{english}{tracking} fiable. La evidencia se acumula reproyectando máscaras
sobre los puntos, así que si el mapa deriva o el \foreignlanguage{english}{tracking} falla,
esa reproyección cae donde no toca y genera disputas que no responden a ningún conflicto
real: la historia de un punto deja de ser comparable consigo misma y el razonamiento se
rompe. Pero esto es, de nuevo, una responsabilidad ajena. Mantener la geometría sana bajo
los puntos es tarea del \foreignlanguage{english}{refresh} geométrico y los
\foreignlanguage{english}{submaps} (\ref{ch:geometric_consistency}); el contest da esa
condición por resuelta y se ocupa solo de lo suyo, arbitrar la disputa.

\section{The contest mechanism}
\label{sec:contest_mechanism}

El contest es el mecanismo que asume la responsabilidad que faltaba: registra los conflictos
entre instancias, acumula la evidencia de cada uno y, a partir de esa evidencia, decide qué
hacer. Todo él se organiza en torno a un par de roles. Cuando dos instancias se disputan un
punto, la que lo posee actualmente es el \emph{defender} $D$ y la que se lo reclama es el
\emph{challenger} $C$. La unidad sobre la que el mecanismo razona y decide no es una
instancia suelta, sino el par ordenado $(D, C)$: toda la maquinaria (la evidencia que
almacena, las señales que calcula, la decisión que emite) está referida a ese par. Su salida
es un \emph{veredicto} por par, $v(D, C) \in \{\textsc{merge}, \textsc{split},
\textsc{no-action}\}$, que traduce la evidencia acumulada en una acción concreta sobre el
mapa o en la ausencia de ella.

El mecanismo se descompone en tres etapas, cada una responsable de una cosa y ninguna más.
El \emph{registro} (\ref{sec:contest_registry}) anota la disputa punto a punto y la agrega
por par, sin decidir. El \emph{discriminador} (\ref{sec:contest_discriminator}) lee esa
evidencia y emite el veredicto, sin tocar el mapa. El \emph{actuador}
(\ref{sec:contest_actuator}) ejecuta el veredicto sobre el mapa. Registro y discriminador
son puros, y esa pureza es deliberada: la disputa se contabiliza y se juzga por separado de
donde se aplica, y solo la última etapa modifica nada (Fig.~\ref{fig:contest_pipeline}).

\begin{figure}[ht]
  \centering
  % PLACEHOLDER: diagrama del pipeline del mecanismo, a dibujar (TikZ o drawio -> PDF).
  % Tres cajas en cadena, con inputs/outputs explícitos entre ellas:
  %   [pasos de tracking: máscaras SAM reproyectadas]
  %      -> REGISTRO -> (señales por par: containment, confidence, exclusivity, ...)
  %      -> DISCRIMINADOR -> (veredicto v(D,C) in {merge, split, no-action})
  %      -> ACTUADOR -> (mapa mutado: instancias fusionadas / puntos reasignados)
  % Marcar registro y discriminador como PUROS (no tocan el mapa); solo el actuador muta.
  \fbox{\parbox[c][3.2cm][c]{0.9\linewidth}{\centering\footnotesize\itshape
  Pipeline del contest: tracking $\to$ registro (señales por par) $\to$ discriminador
  (veredicto) $\to$ actuador (mapa mutado). Registro y discriminador puros; solo el actuador
  toca el mapa.}}
  \caption{Las tres etapas del contest y el flujo de datos entre ellas.}
  \label{fig:contest_pipeline}
\end{figure}

\subsection{The registry}
\label{sec:contest_registry}

El registro es \emph{online}: se construye a la vez que el mapa, sin pasada posterior. En
cada paso de \foreignlanguage{english}{tracking}, cada máscara de SAM se reproyecta sobre los
puntos visibles y resuelve a una instancia, la mayoritaria entre los que ya tienen dueño bajo
la máscara. Sea $p$ un punto de dueño $D = o(p)$; si cae bajo una máscara que resuelve a
$C \neq D$, hay un \emph{robo} (\foreignlanguage{english}{grab}) de $p$ por $C$ contra $D$
(Fig.~\ref{fig:contest_grab}). No se aplica: se anota, sumando un robo a la cuenta de $C$
sobre ese punto.

\begin{figure}[ht]
  \centering
  % PLACEHOLDER: diagrama esquemático del robo, a dibujar (TikZ o drawio -> PDF).
  % Contenido: cámara y máscara SAM reproyectadas sobre la nube de puntos; un punto cuyo
  % dueño actual (defender) difiere de la instancia a la que resuelve la máscara (challenger);
  % flecha que marca el grab que se anota. Opcional: partición leal / robo / huérfano.
  \fbox{\parbox[c][2cm][c]{0.85\linewidth}{\centering\footnotesize\itshape
  Diagrama del robo: máscara SAM reproyectada sobre la nube; un punto cuyo defender difiere
  del challenger de la máscara; se anota el grab.}}
  \caption{El robo: la observación atómica del contest.}
  \label{fig:contest_grab}
\end{figure}

El registro tiene dos niveles. El primero es \emph{por punto}: tres contadores indexados por
el id permanente de cada punto (Tabla~\ref{tab:contest_primary}), la verdad de base que gotea
paso a paso. Cada punto lleva su cuenta de robos (por challenger), de reclamos y de
avistamientos, y nada más se guarda a este nivel; todo lo demás se deriva de aquí.

\begin{table}[ht]
  \centering
  \small
  \begin{threeparttable}
  \renewcommand{\arraystretch}{1.5}
  \begin{tabular}{@{}l l p{7.4cm}@{}}
    \toprule
    \textbf{Signal} & \textbf{Store} & \textbf{Steps counted} \\
    \midrule
    \textbf{\texttt{grabs}}     & $\texttt{point\_id}\mapsto\{C\!:\#\mathrm{KF}\tnote{*}\}$ &
      The point, owned by $D$, fell under $C$'s mask (a grab) \\
    \textbf{\texttt{claims}}    & $\texttt{point\_id}\mapsto\#\mathrm{KF}$ &
      The point fell under any used mask ($C$ or $D$) \\
    \textbf{\texttt{sightings}} & $\texttt{point\_id}\mapsto\#\mathrm{KF}$ &
      The point reprojected into view, with or without a mask \\
    \bottomrule
  \end{tabular}
  \begin{tablenotes}
    \footnotesize\item[*]\itshape $\#\mathrm{KF}$: number of
      \foreignlanguage{english}{keyframe} steps.
  \end{tablenotes}
  \end{threeparttable}
  \caption{Primary signals: per-point counters accumulated in the store.}
  \label{tab:contest_primary}
\end{table}

El segundo nivel agrega esos contadores \emph{por par} $(D, C)$, la unidad de la decisión.
Pero antes de agregar hay que separar la señal del ruido, y de eso se ocupa la
\textbf{firmeza} (\foreignlanguage{english}{firmness}): un robo de un punto cuenta como firme
solo si el challenger se lo llevó en suficientes pasos (\emph{evidencia}, un mínimo de robos)
y esos robos son una fracción real de los reclamos del punto (\emph{compromiso}, la razón
robos/reclamos por encima de un umbral). Un roce de refilón, o un parpadeo de pocos pasos, no
es firme. El número de puntos firmes de un par es \texttt{firm\_points}, y sobre ese
subconjunto se construyen todas las señales secundarias (Tabla~\ref{tab:contest_secondary}).
Estas viven cacheadas por defender: un robo solo ensucia a su defender, cuyas señales se
recomputan en el siguiente refresco, sin re-derivar el resto.

\begin{table}[ht]
  \centering
  \small
  \begin{threeparttable}
  \renewcommand{\arraystretch}{1.5}
  \begin{tabular}{@{}l p{11.6cm}@{}}
    \toprule
    \textbf{Signal} & \textbf{How it is computed} \\
    \midrule
    \textbf{\texttt{firm\_points}} & points of $D$ that $C$ steals firmly \\
    \textbf{\texttt{total\_grabs}} & sum of all $C$'s grabs over $D$'s points (firm or not) \\
    \textbf{\texttt{containment}}  & $\mathrm{firm\_points}\,/\,|D|\tnote{*}$ \\
    \textbf{\texttt{confidence}}   & per-point $\mathrm{grabs}_C/\mathrm{claims}$, averaged over the disputed points \\
    \textbf{\texttt{exclusivity}}  & firm points disputed only by $C$ $\,/\,\mathrm{firm\_points}$ \\
    \bottomrule
  \end{tabular}
  \begin{tablenotes}
    \footnotesize\item[*]\itshape $|D|$: number of points of instance $D$.
  \end{tablenotes}
  \end{threeparttable}
  \caption{Secondary signals: per-pair $(D,C)$ aggregates over firm points.}
  \label{tab:contest_secondary}
\end{table}

Este es todo el registro del contest: un flujo de robos por punto que gotea en cada paso,
tres contadores primarios por punto y un agregado por par vivo por defender. En ningún
momento se ha tocado el mapa ni se ha revocado una asignación; el registro es una
contabilidad fiel de la disputa, deliberadamente separada de la decisión. Convertir esa
evidencia en un veredicto es competencia del discriminador.

\subsection{The discriminator}
\label{sec:contest_discriminator}

El discriminador toma los pares $(D, C)$ de un defender, con sus señales ya calculadas, y
emite el veredicto. Es una función pura: no mira el mapa ni el descriptor, solo la evidencia
que el registro acumuló. Toda ella se ordena en torno a una única cantidad, la contención
$\mathrm{containment}(D, C)$: qué fracción del defender le roba el challenger con firmeza. La
contención decide la \emph{forma} de la acción, y la \foreignlanguage{english}{confidence} y
la exclusividad deciden
si hay \emph{evidencia suficiente} para ejecutarla.

Antes de razonar nada, un filtro de masa descarta el ruido: un par con menos de
$\tau_{\text{masa}}$ robos crudos ($\mathrm{total\_grabs} < 50$) no es una disputa, es un
roce, y se ignora. Sobre los pares que sobreviven, la contención reparte la decisión en tres
bandas.

\begin{enumerate}
  \item \textbf{Contención alta} ($\mathrm{containment} \geq 0.6$): el challenger se lleva casi
  todo el defender, luego el defender es un \emph{fragmento} suyo (casos A y C). La acción
  natural es \textsc{merge}, pero solo si hay un dueño fiable: se exige una
  \foreignlanguage{english}{confidence} $\mathrm{confidence} \geq 0.5$ (de media, cada punto se roba
  en la mayoría de sus reclamos) y que
  todos los challengers firmes sean, vía la fusión, una sola instancia. Si lo contienen
  varios objetos distintos no hay dueño claro y se emite \textsc{no-action}.

  \item \textbf{Contención parcial} ($0.4 \leq \mathrm{containment} < 0.6$): el challenger se lleva
  una parte grande, pero no tanta como para que la sola magnitud baste. Aquí el umbral sube:
  \textsc{merge} solo si el trozo es \emph{limpio}, es decir con confidence alta
  ($\mathrm{confidence} \geq 0.6$) y exclusivo ($\mathrm{exclusivity} \geq 0.8$, casi todos sus puntos los
  disputa este challenger y nadie más), y de nuevo con raíz única. La evidencia de magnitud
  es más débil, así que se compensa exigiendo más limpieza en las otras dos señales.

  \item \textbf{Contención baja} ($\mathrm{containment} < 0.4$): el defender \emph{retiene} la
  mayoría de sus puntos, pero cede trozos exclusivos a uno o varios vecinos. No es un
  fragmento de nadie, es un objeto que arrastra pedazos de otros (caso B). La acción es
  \textsc{split}, y no una sino tantas como haga falta: por cada challenger con un trozo
  con confidence alta ($\mathrm{confidence} \geq 0.6$) y muy exclusivo ($\mathrm{exclusivity} \geq 0.9$) se
  desprende ese trozo hacia él. Si ningún challenger cualifica, el par está en el borde y se
  emite \textsc{no-action}.
\end{enumerate}

El umbral de exclusividad sube de $0.8$ en el \textsc{merge} parcial a $0.9$ en el
\textsc{split}, y no por capricho: un split parte geometría en dos y es más caro de deshacer
si sobra, así que se exige que la frontera del trozo sea más nítida antes de cortar. El suelo
$\mathrm{containment} = 0.4$ que separa split de merge es empírico, fijado sobre el conjunto
de \foreignlanguage{english}{tuning} (\ref{sec:contest_results}).

Visto en conjunto, el \textsc{merge} clásico reaparece como el extremo de esta escala: el
caso en que la contención es total y el defender migra entero. El discriminador solo extiende
esa misma idea a las contenciones intermedias y bajas, donde la decisión deja de ser mover
toda la instancia y pasa a mover exactamente los puntos que la evidencia respalda.

\subsection{The actuator}
\label{sec:contest_actuator}
% Ejecuta el veredicto sobre el mapa: MERGE (union-find de instancias) y SPLIT (mueve los
% split_points al challenger). Última etapa, la única que muta. Aquí se cuela el timing/coste
% real (Real-time ya no es sección propia): coste por refresco, cuello de botella y rediseño
% online. Doc: docs/contest_timing_analysis.html

\section{Results}
\label{sec:contest_results}

Los umbrales del discriminador son empíricos, así que todos los números se reportan sobre un
conjunto de escenas que el ajuste nunca vio. Las ocho escenas de Replica se parten en
\foreignlanguage{english}{tuning} (office0, office3, office4, room2), las únicas sobre las que
se fijan los umbrales, y \foreignlanguage{english}{held-out} (office1, office2, room0, room1),
que se mantienen cerradas durante todo el ajuste y se abren una sola vez, al final. Ningún
umbral se toca mirando el \foreignlanguage{english}{held-out}; el $\mathrm{AP}$ que se defiende
es siempre el suyo, de modo que mide generalización y no amoldamiento.

% PRELIMINAR: números a recalcular antes de la entrega (placeholder para ver la maqueta).
Sobre la trayectoria limpia de Replica, sin saltos inyectados, el contest sube el
$\mathrm{AP}_{\text{agnostic}}$ global de $0.234$ a $0.249$ ($+0.015$, $+6.4\%$) respecto al
mapa crudo (\emph{baseline} sin fusión sobre GT simulado). La mejora aparece en las ocho
escenas, y se sostiene en el \foreignlanguage{english}{held-out}: la ganancia media es de
$+9.7\%$ en las escenas de \foreignlanguage{english}{tuning} y de $+4.55\%$ en las de
\foreignlanguage{english}{held-out}, positiva en ambos casos
(Tabla~\ref{tab:contest_ap_scenes}).

\begin{table}[ht]
  \centering
  \small
  \renewcommand{\arraystretch}{1.2}
  \begin{tabular}{@{}l r r r r@{}}
    \toprule
    \textbf{Scene} & \textbf{Baseline} & \textbf{Contest} & \textbf{$\Delta$ abs} &
      \textbf{$\Delta$ \%} \\
    \midrule
    \multicolumn{5}{@{}l}{\textbf{Tuning}} \\
    office0 & 0.199 & 0.220 & $+0.021$ & $+10.6\%$ \\
    office3 & 0.204 & 0.210 & $+0.006$ & $+2.9\%$  \\
    office4 & 0.236 & 0.280 & $+0.044$ & $+18.6\%$ \\
    room2   & 0.232 & 0.247 & $+0.015$ & $+6.5\%$  \\
    \addlinespace[1pt]
    \textbf{mean} & \textbf{0.218} & \textbf{0.239} & $\mathbf{+0.021}$ & $\mathbf{+9.7\%}$ \\
    \midrule
    \multicolumn{5}{@{}l}{\textbf{Held-out}} \\
    office1 & 0.170 & 0.175 & $+0.005$ & $+2.9\%$  \\
    office2 & 0.210 & 0.225 & $+0.015$ & $+7.1\%$  \\
    room0   & 0.269 & 0.287 & $+0.018$ & $+6.7\%$  \\
    room1   & 0.344 & 0.349 & $+0.005$ & $+1.5\%$  \\
    \addlinespace[1pt]
    \textbf{mean} & \textbf{0.248} & \textbf{0.259} & $\mathbf{+0.011}$ & $\mathbf{+4.55\%}$ \\
    \midrule
    \textbf{All} & \textbf{0.234} & \textbf{0.249} & $\mathbf{+0.015}$ & $\mathbf{+6.4\%}$ \\
    \bottomrule
  \end{tabular}
  \caption{$\mathrm{AP}_{\text{agnostic}}$ per Replica scene, baseline vs.\ contest.}
  \label{tab:contest_ap_scenes}
\end{table}

Las métricas \emph{class-aware} (mIoU / mAcc) no empeoran: se mueven ligeramente al alza
(mIoU $+0.2$, mAcc $+0.3$ globales, Tabla~\ref{tab:contest_semantic}). Es coherente con la
naturaleza del mecanismo: el contest reorganiza la propiedad de los puntos entre instancias,
no cambia la etiqueta semántica que cada una recibe, así que su efecto vive en la granularidad
de las instancias (lo que mide el $\mathrm{AP}$ agnóstico) y deja el mIoU por clase
prácticamente intacto.

\begin{table}[ht]
  \centering
  \small
  \renewcommand{\arraystretch}{1.2}
  \begin{tabular}{@{}l r r r r r r@{}}
    \toprule
    & \multicolumn{3}{c}{\textbf{mIoU}} & \multicolumn{3}{c}{\textbf{mAcc}} \\
    \cmidrule(lr){2-4} \cmidrule(lr){5-7}
    \textbf{Scene} & Base & Contest & $\Delta$ & Base & Contest & $\Delta$ \\
    \midrule
    office0 & 25.4 & 25.9 & $+0.5$ & 33.9 & 34.3 & $+0.4$ \\
    office1 & 16.2 & 16.1 & $-0.1$ & 31.7 & 31.7 & $+0.0$ \\
    office2 & 24.6 & 25.9 & $+1.3$ & 32.4 & 34.0 & $+1.5$ \\
    office3 & 25.8 & 25.9 & $+0.0$ & 32.3 & 32.3 & $+0.0$ \\
    office4 & 43.1 & 42.7 & $-0.3$ & 55.4 & 55.8 & $+0.4$ \\
    room0   & 33.2 & 33.6 & $+0.4$ & 39.9 & 40.3 & $+0.4$ \\
    room1   & 38.6 & 40.2 & $+1.6$ & 53.8 & 56.1 & $+2.3$ \\
    room2   & 26.6 & 26.6 & $+0.0$ & 35.1 & 35.1 & $+0.0$ \\
    \midrule
    \textbf{All} & \textbf{27.1} & \textbf{27.3} & $\mathbf{+0.2}$ & \textbf{40.2} & \textbf{40.5} & $\mathbf{+0.3}$ \\
    \bottomrule
  \end{tabular}
  \caption{Class-aware semantic metrics (\%) per Replica scene, baseline vs.\ contest.}
  \label{tab:contest_semantic}
\end{table}
% Validación en AISLAMIENTO, sobre trayectoria sana y sin saltos: se atribuye la mejora
% al mecanismo, no a su interacción con el ruido. COROLARIO, que hay que escribir: el
% contest PRESUPONE geometría sana y tracking fiable -> el sustrato no es un preámbulo,
% es una PRECONDICIÓN.
%
% ABLACIÓN 2x2 (el experimento estrella de la memoria), sobre saltos inyectados:
%
%                    | sin submapas            | con submapas
%   .................+.........................+...........................
%   sin contest      | baseline degradado      | ¿el submapa solo, ayuda?
%   con contest      | contest sobre sustrato  | la propuesta completa
%                    | roto -> debería ir mal  |
%
% La celda de abajo-izquierda es la que DEMUESTRA que el orden causal es real y no
% retórica. Asegurar este experimento sí o sí.
